Condition throughout on \(\mathcal T\), and fix an arm \(a\). Write \(\mathcal O_a=\{O\in\mathcal O:O\subseteq\mathcal U_a\}\). By \(\texttt{lem:orbit-score-factorization}\), a calibration score and the target score have conditional laws
\[
\mathsf Q_a=\sum_{O\in\mathcal O_a}q_{a,O}F_{a,O},
\qquad
\mathsf P_a=\sum_{O\in\mathcal O_a}p_{a,O}F_{a,O},
\tag{1}
\]
and the \(M\) calibration scores are independent and identically distributed from \(\mathsf Q_a\), independently of the target score from \(\mathsf P_a\).

For a measurable \(f:[0,\infty)\to[0,1]\), set \(h_O=\int f\,dF_{a,O}\) and \(d_O=p_{a,O}-q_{a,O}\). Since both orbit-mass vectors are probability vectors, \(\sum_Od_O=0\). The positive and negative parts of \((d_O)_O\) therefore have the same total mass, namely
\[
\frac12\sum_{O\in\mathcal O_a}|d_O|=\Delta_a(w).
\]
It follows that
\[
\left|\int f\,d(\mathsf P_a-\mathsf Q_a)\right|
=\left|\sum_{O\in\mathcal O_a}d_Oh_O\right|
\le\Delta_a(w).
\tag{2}
\]
Thus passage through the common orbit-specific score kernels contracts total variation.

On an auxiliary product extension introduce \(S'_{0,a}\sim\mathsf Q_a\), independent of the calibration scores. Append independent continuous tie-breakers as prescribed by \(\texttt{def:rank-inversion}\). The augmented rank of \(S'_{0,a}\) among the \(N=M+1\) independent and identically distributed augmented scores is uniform on \(\{1,\ldots,N\}\). If the unaugmented comparison score strictly exceeds the reported \(k_a\)-th calibration order statistic, its augmented rank is greater than \(k_a\), regardless of the tie-breaker. Hence
\[
\Pr\{S'_{0,a}>r_{a,(k_a)}\mid\mathcal T\}
\le\frac{N-k_a}{N}=\rho_{M,\varepsilon_a}.
\tag{3}
\]
This also covers \(k_a=N\), when the event is empty by the infinite-threshold convention.

For each fixed calibration realization, apply \((2)\) to \(f(s)=\mathbf 1\{s>r_{a,(k_a)}\}\). The actual target is independent of the calibration sample and has law \(\mathsf P_a\), so averaging the resulting comparison and using \((3)\) yields
\[
\Pr_P\{S_{0,a}>r_{a,(k_a)}\mid\mathcal T\}
\le\rho_{M,\varepsilon_a}+\Delta_a(w).
\tag{4}
\]
Let \(\widetilde I^{\boldsymbol\varepsilon}_{a,w}\) be the un-clipped residual interval. On \(\{S_{0,a}\le r_{a,(k_a)}\}\), the target outcome lies in that interval. The bounded-outcome assumption also gives \(Y_{0T}(a)\in[-K,K]\). Since
\[
I^{\boldsymbol\varepsilon}_{a,w}=\widetilde I^{\boldsymbol\varepsilon}_{a,w}\cap[-K,K],
\]
we have
\[
\{Y_{0T}(a)\notin I^{\boldsymbol\varepsilon}_{a,w}\}
\subseteq
\{S_{0,a}>r_{a,(k_a)}\}.
\tag{5}
\]
If neither arm fails in \((5)\), then
\[
\tau_T=Y_{0T}(1)-Y_{0T}(0)
\in I^{\boldsymbol\varepsilon}_{1,w}-I^{\boldsymbol\varepsilon}_{0,w}
=\mathcal C^{\boldsymbol\varepsilon}_w.
\tag{6}
\]
The union bound applied to \((4)\)--\((6)\), without any cross-arm independence, proves the displayed coverage lower bound after truncation at zero. The coordinatewise elbow description follows from \(\texttt{lem:armwise-finite-rank-staircase}\).

Finally, \(\sum_O\delta_{a,O}(w)=0\). Hence the positive-coordinate sum and the absolute negative-coordinate sum both equal \(\Delta_a(w)\), and every coordinate satisfies \(|\delta_{a,O}(w)|\le\Delta_a(w)\). The diagonal specialization follows from the last assertion of \(\texttt{lem:armwise-finite-rank-staircase}\) and the symmetric identities in \(\texttt{def:rank-inversion}\).